\BourbakiExercisesSection

\begin{enumerate}
\def\labelenumi{\arabic{enumi}.}
\item
  Let N be the Z-module Z/4Z, M the submodule 2Z/4Z of N and \(j:M\to N\) the canonical injection. Show that \(\mathsf{T}(j):\mathsf{T}(\mathsf{M})\to\mathsf{T}(\mathsf{N})\) is not injective.
\item
  Let I and J be two finite sets, \(I_{0} = I \cup \{\alpha\}\) , \(J_{0} = J \cup \{\beta\}\) , where \(\alpha\) and \(\beta\) are two elements not belonging to I and J respectively, and suppose that \(I_{0}\) and \(J_{0}\) are disjoint. Let E be a vector space of finite rank over a commutative field K and let z be an element of \(\mathbf{T}_{J}^{\mathrm{I}}(\mathbf{E})\) . For all \(i \in I\) , let \(W_{i}'\) be the subspace of \(E^{*}\) consisting of the \(x^{*} \in E^{*}\) such that \(c_{\beta}^{i}(z \otimes x^{*}) = 0\) ( \(c_{\beta}^{i}\) being the contraction of index \(i \in I\) and index \(\beta \in J_{0}\) in \(\mathbf{T}_{J_{0}}^{\mathrm{I}}(\mathbf{E})\) ) and let \(V_{i}\) be the subspace of E orthogonal to \(W_{i}'\) . For all \(j \in J\) , let \(W_{j}\) be the subspace of E consisting of the \(x \in E\) such that \(c_{j}^{\alpha}(z \otimes x) = 0\) ( \(c_{j}^{\alpha}\) being the contraction of index \(\alpha \in I_{0}\) and index \(j \in J\) in \(T_{J}^{I_{0}}(E)\) ) and let \(V_{j}'\) be the subspace of \(E^{*}\) orthogonal to \(W_{j}\) . Show that z belongs to the tensor product \(\left(\bigotimes_{i \in I} V_{i}\right) \otimes \left(\bigotimes_{j \in J} V_{j}'\right)\) canonically identified with a subspace of \(T_{J}^{I}(E)\) . Moreover, if \((\mathbf{U}_{i})_{i \in I}\) is a family of subspaces of E, \((\mathbf{U}_{j}')_{j \in J}\) a family of subspaces of \(E^{*}\) such that z belongs to the tensor product \(\left(\bigotimes_{i \in I} U_{i}\right) \otimes \left(\bigotimes_{j \in J} U_{j}'\right)\) , identified with a subspace of \(T_{J}^{I}(E)\) , then \(V_{i} \subset U_{i}\) and \(V_{j}' \subset U_{j}'\) for all \(i \in I\) and \(j \in J\) (cf.~II, §7, no.8, Proposition 17).
\item
  Let M be a finitely generated projective A-module. For every endomorphism u of M, let \(\tilde{u}\) denote the tensor \(\theta_{\mathbf{M}}^{-1}(u) \in \mathbf{M}^{*} \otimes_{\mathbf{A}} \mathbf{M}\) corresponding to it.
\end{enumerate}

\begin{enumerate}
\def\labelenumi{(\alph{enumi})}
\item
  For every element \(x \in \mathbf{M}\) , show that \(u(x) = c_2^1(x \otimes \tilde{u})\) , where \(x \otimes \tilde{u}\) belongs to \(\mathbf{M} \otimes \mathbf{M}^* \otimes \mathbf{M} = \mathsf{T}_{\langle 2\rangle}^{(1,3)}(\mathbf{M})\) .
\item
  Let \(u, v\) be two endomorphisms of \(M\) and \(w = u \circ v\) ; show that \(\tilde{w} = c_3^2(\tilde{v} \otimes \tilde{u})\) , where \(\tilde{v} \otimes \tilde{u}\) belongs to \(T_{(1,3)}^{(2,4)}(M)\) .
\item
  For every automorphism \(s\) of \(\mathbf{M}\) and every tensor \(z \in \mathbf{T}_{\mathbf{J}}^{\mathbf{I}}(\mathbf{M})\) , a tensor \(s, z \in \mathbf{T}_{\mathbf{J}}^{\mathbf{I}}(\mathbf{M})\) is defined by the condition that if \(z = \bigotimes_{k \in \mathbf{I} \cup \mathbf{J}} z_k\) where \(z_k \in \mathbf{M}\) for \(k \in \mathbf{I}\) and \(z_k \in \mathbf{M}^*\) for \(k \in \mathbf{J}\) , then \(s, z = \bigotimes_{k \in \mathbf{I} \cup \mathbf{J}} z_k'\) with \(z_k' = s(z_k)\) if \(k \in \mathbf{I}\) , \(z_{k}^{\prime}=^{t}s^{-1}(z_{k})\) if \(k\in J\) . Show that the group \(\mathbf{GL}(\mathbf{M})\) operates by the law \((s,z)\mapsto s.z\) on the A-module \(T_{J}^{I}(M)\) , the mappings \(z\mapsto s.z\) being automorphisms of this A-module. Moreover, for every ordered pair of indices \(i\in I\) , \(j\in J\) ,
\end{enumerate}

\[
c _ {j} ^ {i} (s. z) = s. c _ {j} ^ {i} (z).
\]

If M is projective and finitely generated, show, in the notation of Exercise 2, that \(s. \tilde{u} = (s \circ u \circ s^{-1})^{\sim}\) .

\begin{enumerate}
\def\labelenumi{\arabic{enumi}.}
\setcounter{enumi}{3}
\tightlist
\item
  Let M be a finitely generated projective A-module; for every bilinear mapping \(u\) of \(\mathbf{M} \times \mathbf{M}\) into \(\mathbf{M}\) , let \(\tilde{u}\) denote the tensor \(\theta_{\mathbf{M}}^{-1}(u)\) in \(\mathsf{T}_{(1,2)}^{(3)}(\mathbf{M})\) . Show that for the algebra structure on \(\mathbf{M}\) defined by \(u\) to be associative, it is necessary and sufficient that the tensors \(c_4^3 (\tilde{u} \otimes \tilde{u})\) and \(c_2^6 (\tilde{u} \otimes \tilde{u})\) correspond under the canonical associativity isomorphism of \(\mathsf{T}_{(1,2,3)}^{(4)}(\mathbf{M})\) onto \(\mathsf{T}_{(1,3,4)}^{(2)}(\mathbf{M})\) .
\end{enumerate}

¶ 5. Let E be a vector space over a commutative field K and let T(E) be the tensor algebra of E.

\begin{enumerate}
\def\labelenumi{(\alph{enumi})}
\item
  Show that \(\mathsf{T}(\mathbf{E})\) contains no divisors of 0 and that the only invertible elements in \(\mathsf{T}(\mathbf{E})\) are the scalars \(\neq 0\) . If \(\dim (\mathbf{E}) > 1\) , \(\mathsf{T}(\mathbf{E})\) is a non-commutative K-algebra.
\item
  Let u, v be two elements of T(E). Show that, if there exist two elements a, b of T(E) such that \(ua = vb \neq 0\) , one of the elements u, v is a right multiple of the other. (Consider first the case where u and v are homogeneous; reduce it to the case where E is finite-dimensional and write \(u = e_{1}u_{1} + \cdots + e_{n}u_{n}\) , \(v = e_{1}v_{1} + \cdots + e_{n}v_{n}\) , where \((e_{k})\) is a basis of E. Conclude by arguing by induction on the smallest degree of the homogeneous components of ua.) Deduce that, if \(\dim(\mathrm{E}) > 1\) , the ring T(E) admits no field of left fractions (I, § 9, Exercise 15).
\item
  Show that for two non-zero elements \(u, v\) of \(\mathsf{T}(\mathbf{E})\) to be permutable, it is necessary and sufficient that there exist a vector \(x \neq 0\) in \(\mathbf{E}\) such that \(u\) and \(v\) belong to the subalgebra of \(\mathsf{T}(\mathbf{E})\) generated by \(x\) (use (b)). Deduce that, if \(\dim(\mathbf{E}) > 1\) , the centre of \(\mathsf{T}(\mathbf{E})\) is equal to \(\mathbf{K}\) .
\end{enumerate}

¶ 6. Let K be a commutative field and E, F two K-algebras each with a unit element which is identified with the unit element l of K. Consider the K-algebra T(E ⊕ F), where E ⊕ F is considered as a vector space over K; E (resp. F) is identified with a vector subspace of T(E ⊕ F) under the canonical injection. Let \(\mathfrak{J}\) be the two-sided ideal of T(E ⊕ F) generated by the elements \(x_1 \otimes x_2 - (x_1 x_2)\) and \(y_1 \otimes y_2 - (y_1 y_2)\) , where \(x_i \in E\) , \(y_i \in F\) , \(i = 1, 2\) ; let E * F denote the quotient algebra T(E ⊕ F)/\(\mathfrak{J}\); it is called the free product of the K-algebras E and F; let \(\phi\) denote the canonical homomorphism of T(E ⊕ F) onto E * F and \(\alpha\) and \(\beta\) the restrictions of \(\phi\) to E and F; \(\alpha\) and \(\beta\) are K-algebra homomorphisms.

\begin{enumerate}
\def\labelenumi{(\alph{enumi})}
\item
  Show that for every ordered pair of K-homomorphisms \(u: \mathbf{E} \to \mathbf{G}\) , \(v: \mathbf{F} \to \mathbf{G}\) into a K-algebra \(\mathbf{G}\) , there exists one and only one K-homomorphism \(w: \mathbf{E} * \mathbf{F} \to \mathbf{G}\) such that \(u = w \circ \alpha\) and \(v = w \circ \beta\) (which justifies the terminology).
\item
  Let \(R_{n}\) be the vector subspace of \(\mathsf{T}(\mathbf{E}\oplus\mathbf{F})\) the sum of the \(\mathsf{T}^{k}(\mathbf{E}\oplus\mathbf{F})\) for \(0\leqslant k\leqslant n\) and write \(J_{n}=J\cap R_{n}\) . Show that, if we consider a basis of E consisting of l and a family \((e_{\lambda})_{\lambda\in\mathbb{L}}\) and a basis of F consisting of l and a family \((f_{\mu})_{\mu\in\mathbb{M}}\) , then a supplementary subspace of \(J_{n}+R_{n-1}\) is obtained in \(R_{n}\) by considering the subspace of \(R_{n}\) with basis the elements of the form \(z_{1}\otimes z_{2}\otimes\cdots\otimes z_{n}\) where, either \(z_{2j+1}\) is equal to one of the \(e_{\lambda}\) for \(2j+1\leqslant n\) and \(z_{2j}\) is equal to one of the \(f_{\mu}\) for \(2j\leqslant n\) , or \(z_{2j+1}\) is equal to one of the \(f_{\mu}\) for \(2j+1\leqslant n\) and \(z_{2j}\) is equal to one of the \(e_{\lambda}\) for \(2j\leqslant n\) (argue by induction on n).
\item
  Let \(\mathbf{P}_n = \phi (\mathbf{R}_n)\subset \mathbf{E}*\mathbf{F}\) ; then \(\mathbf{P}_n\subset \mathbf{P}_{n + 1},\mathbf{P}_0 = \mathbf{K},\mathbf{P}_h\mathbf{P}_k\subset \mathbf{P}_{h + k}\) and \(\mathbf{E}*\mathbf{F}\) is the union of the \(\mathbf{P}_n\) . Show that there is a canonical vector space isomorphism, for \(n\geqslant 1\)
\end{enumerate}

\[
\mathrm{P} _ {n} / \mathrm{P} _ {n - 1} \rightarrow \left(\mathrm{G} _ {1} \otimes \mathrm{G} _ {2} \otimes \dots \otimes \mathrm{G} _ {n}\right) \oplus \left(\mathrm{G} _ {2} \otimes \mathrm{G} _ {3} \otimes \dots \otimes \mathrm{G} _ {n + 1}\right)
\]

where we write \(G_{k} = E/K\) for k odd, \(G_{k} = F/K\) for k even. In particular the homomorphisms \(\alpha\) and \(\beta\) are injective, which allows us to identify E and F with sub-K-algebras of E * F. \(P_{n}^{\prime}\) (resp. \(P_{n}^{\prime\prime}\) ) is used to denote the inverse images in P of \(G_{1} \otimes G_{2} \otimes \cdots \otimes G_{n}\) (resp. \(G_{2} \otimes G_{3} \otimes \cdots \otimes G_{n+1}\) ); the height of an element \(z \in E * F\) , denoted by \(h(z)\) , is the smallest integer such that \(z \in P_{n}\) ; if \(h(z) = n\) and \(z \in P_{n}^{\prime}\) (resp. \(z \in P_{n}^{\prime\prime}\) ), z is called pure odd (resp. pure even). Then \(P_{n} = P_{n}^{\prime} + P_{n}^{\prime\prime}\) and \(P_{n}^{\prime} \cap P_{n}^{\prime\prime} = P_{n-1}\) for \(n \geqslant 1\) .

\begin{enumerate}
\def\labelenumi{(\alph{enumi})}
\setcounter{enumi}{3}
\item
  Show that, if \(u, v\) are two elements of \(\mathbf{E} * \mathbf{F}\) such that \(h(u) = r\) , \(h(v) = s\) , then \(h(uv) \leqslant h(u) + h(v)\) and that \(h(uv) = h(u) + h(v)\) unless \(u\) and \(v\) are pure and, either \(r\) is even and \(u, v\) are not of the same parity, or \(r\) is odd and \(u, v\) are of the same parity. (If for example \(r\) is even, \(u \in \mathbf{P}_r'\) , \(v \in \mathbf{P}_s'\) , show that \(uv \in \mathbf{P}_{r+s-1}\) is possible only if \(u \in \mathbf{P}_{r-1}\) or \(v \in \mathbf{P}_{s-1}\) .)
\item
  Suppose that E has no divisors of zero and, in the notation of (d), suppose that r is even, u pure even and v pure odd. Show that \(h(uv) = r + s - 1\) , unless there exist two invertible elements x, y of E such that \(uy \in P_{r-1}\) and \(xb \in P_{s-1}\) . (Let \((a_i)\) (resp. \((b_j)\) ) be the basis of the supplementary subspace of \(J_r + R_{r-1}\) in \(R_r\) (resp. of \(J_s + R_{s-1}\) in \(R_s\) ) considered in (b); write \(u = \sum_i \phi(a_i)x_i\) and \(v = \sum_j y_j \phi(b_j)\) , where \(x_i\) and \(y_j\) are in E, and show that, if \(uv \in P_{r+s-2}\) , then necessarily \(x_i y_j \in K\) .) Treat similarly the other cases where \(h(uv) < h(u) + h(v)\) when E and F are assumed to have no divisors of zero.
\item
  Deduce from (d) and (e) that, when E and F have no divisors of zero, nor does E * F.
\item
  Generalize the above definitions and results to any finite number of unital K-algebras.
\end{enumerate}

\begin{enumerate}
\def\labelenumi{\arabic{enumi}.}
\setcounter{enumi}{6}
\tightlist
\item
  Let E be a vector space of finite dimension n over a field K. For every integer \(r \geqslant 1\) , the symmetric group \(S_{r}\) operates linearly on the tensor product \(E^{\otimes r}\) by the action \((\sigma, z) \mapsto \sigma \cdot z\) such that
\end{enumerate}

\[
\sigma . \left(x _ {1} \otimes x _ {2} \otimes \dots \otimes x _ {r}\right) = x _ {\sigma^ {- 1} (1)} \otimes x _ {\sigma^ {- 1} (2)} \otimes \dots \otimes x _ {\sigma^ {- 1} (r)}
\]

for every sequence of \(r\) vectors \((x_i)_{1 \leqslant i \leqslant r}\) of \(\mathbf{E}\) . Let \(u_{\sigma}(z) = \sigma, z\) . Show that for every sequence \((v_i)_{1 \leqslant i \leqslant r}\) of \(r\) endomorphisms of \(\mathbf{E}\) ,

\[
\operatorname{Tr} (u _ {\sigma} \circ (v _ {1} \otimes v _ {2} \otimes \dots \otimes v _ {r})) = \operatorname{Tr} (w _ {1}) \operatorname{Tr} (w _ {2}) \dots \operatorname{Tr} (w _ {k})\tag{*}
\]

where the endomorphisms \(w_{j}\) of \(\mathbf{E}\) are defined as follows: \(\sigma^{-1}\) is decomposed into cycles, \(\sigma^{-1} = \lambda_1\lambda_2\ldots \lambda_k\) and, if \(\lambda_{j} = (i_{1}i_{2}\dots i_{h})\) , we write

\[
w _ {j} = v _ {i _ {1}} \circ v _ {i _ {2}} \circ \dots \circ v _ {i _ {h}}.
\]

(Reduce it to calculating the left hand side of (*) when each \(v_{i}\) is of the form \(x \mapsto \langle x, a_{h}^{*} \rangle a_{k}\) , where \((a_{j})_{1 \leqslant j \leqslant n}\) is a basis of E and \((a_{j}^{*})_{1 \leqslant j \leqslant n}\) the dual basis.)

¶ 8. Let A be an integral domain and K its field of fractions; for every A-module M, let τ(M) denote its torsion module (II, § 7, no. 10).

\begin{enumerate}
\def\labelenumi{(\alph{enumi})}
\item
  Show that for the algebra \(\mathsf{T}(\mathbf{M})\) to have no divisors of zero, it is necessary and sufficient that the A-module \(\mathsf{T}(\mathbf{M})\) be torsion-free (observe that \(\mathsf{T}(\mathbf{M})_{(\mathbb{K})}\) is isomorphic to \(\mathsf{T}(\mathbf{M}_{(\mathbb{K})})\) and use Exercise 5).
\item
  Let \(u: \mathbf{M} \to \mathbf{M}'\) be an injective A-module homomorphism. Show that \(\operatorname{Ker}(\mathsf{T}(u)) \subset \tau(\mathsf{T}(\mathbf{M}))\) ; if \(\mathbf{M}'\) is torsion-free, then \(\operatorname{Ker}(\mathsf{T}(u)) = \tau(\mathsf{T}(\mathbf{M}))\) (cf.~II, § 7, no. 10, Proposition 27).
\item
  Give an example of a torsion-free A-module M such that \(\tau(\mathsf{T}(\mathbf{M}))\) is non-zero (II, § 7, Exercise 31).
\end{enumerate}

¶ 9. Let A be a commutative ring, F the free associative algebra over A of the set of integers I = {[}1, n{]}, where n ≥ 2, and X \(_{i}\) (1 ≤ i ≤ n) the corresponding indeterminates, so that a canonical basis of F over A consists of the products X \(_{i_1}\) X \(_{i_2}\) \ldots{} X \(_{i_r}\) , where (i \(_{1}\) , \ldots, i \(_{r}\) ) is an arbitrary finite sequence of elements of I.

\begin{enumerate}
\def\labelenumi{(\alph{enumi})}
\item
  Let \(C_{2}\) denote the set of commutators \([X_{i}, X_{j}] = X_{i}X_{j} - X_{j}X_{i}\) for \(1 \leqslant i < j \leqslant n\) . Assuming that \(C_{m-1}\) is defined, \(C_{m}\) denotes the set of elements \([X_{i}, P] = X_{i}P - PX_{i}\) , where \(1 \leqslant i \leqslant n\) and P runs through \(C_{m-1}\) . Let U be the graded subalgebra of F generated by l and the union of the \(C_{m}\) for \(m \geqslant 2\) . Show that, for \(1 \leqslant i \leqslant n\) , \([X_{i}, P] \in U\) for all \(P \in U\) (confine it to the case where P is a product of elements of \(\bigcup_{m} C_{m}\) and argue by induction on the number of factors).
\item
  Suppose from now on that A is a field. For every integer \(m \geqslant 0\) , let \(\mathbf{U}_m\) be the vector A-space of homogeneous elements of degree \(m\) in U (so that \(\mathbf{U}_0 = \mathbf{A}\) , \(\mathbf{U}_1 = \{0\}\) ); choose a basis \(\mathbb{R}_{m,1}, \ldots, \mathbb{R}_{m,d_m}\) in each \(\mathbf{U}_m\) , consisting of products of elements of \(\bigcup_{k\leqslant m}\mathbf{C}_k\) . For every ordered pair of integers \((m,k)\) such that \(m\geqslant 2,1\leqslant k\leqslant d_m\) , let \(\mathbf{E}_{m,k}\) be the vector sub-A-space of U generated by the \(\mathbb{R}_{p,j}\) such that \(p > m\) or \(p = m\) and \(j\geqslant k\) ; show that \(\mathbf{E}_{m,k}\) is a two-sided ideal of U and that, if \(\mathbf{L}_{m,k}\) is the left ideal of F generated by \(\mathbf{E}_{m,k},\mathbf{L}_{m,k}\) is a two-sided ideal of F.
\item
  For every multiindex \(\alpha = (\alpha_{1},\ldots ,\alpha_{n})\in \mathbf{N}^{n}\) , we write
\end{enumerate}

\[
\mathrm{X} ^ {\alpha} = \mathrm{X} _ {1} ^ {\alpha_ {1}} \mathrm{X} _ {2} ^ {\alpha_ {2}} \dots \mathrm{X} _ {n} ^ {\alpha_ {n}}.
\]

Show that the elements \(X^{\alpha}R_{m,k}\) ( \(m > 0, l \leqslant k \leqslant d_{m}\) ) form a basis of F over K. (To prove that these elements form a generating system of the vector space F, prove that every element \(X_{i_{1}} \ldots X_{i_{r}}\) is a linear combination of them by induction on the degree r. To see that the \(X^{\alpha}R_{m,k}\) are linearly independent, argue by reductio ad absurdum by considering the maximum degree with respect to one of the \(X_{i}\) of the monomials \(X^{\alpha}\) which would appear in a linear relation between the \(X^{\alpha}R_{m,k}\) and proceed by induction on this maximum degree; for this, reduce it to the case where K is infinite (considering if necessary \(F \otimes_{K} K'\) for an extension field \(K'\) of K) and observe that when \(X_{i} + \lambda\) is substituted in an \(R_{m,k}\) for \(X_{i}\) , for some \(\lambda \in K'\) , the element \(R_{m,k}\) is unaltered).

\begin{enumerate}
\def\labelenumi{(\alph{enumi})}
\setcounter{enumi}{3}
\item
  Show that the two-sided ideal \(\mathfrak{J}(\mathbf{F})\) of \(\mathbf{F}\) generated by the commutators \([u,v] = uv - vu\) of elements of \(\mathbf{F}\) is the sum of the \(\mathbf{U}_m\) for \(m\geqslant 2\) ; \(\mathbf{F} / \mathfrak{J}(\mathbf{F})\) is isomorphic to \(\mathrm{A}[\mathrm{X}_1,\dots ,\mathrm{X}_n]\) and hence is an integral domain.
\item
  Show that the algebra U is not a free associative algebra by proving that the quotient \(\mathrm{U}/\Im(\mathrm{U})\) , where \(\Im(\mathrm{U})\) is generated by the commutators of elements of U, admits nilpotent elements \(\neq0\) . (Consider the image \(z'\) in \(\mathrm{U}/\Im(\mathrm{U})\) of the element \(z = [\mathbf{X}_{i}, \mathbf{X}_{j}]\) of U; \(z'^{2} \neq 0\) but \(z'^{4} = 0\) .)
\end{enumerate}

\begin{enumerate}
\def\labelenumi{\arabic{enumi}.}
\setcounter{enumi}{9}
\tightlist
\item
  Let A be a commutative ring,
\end{enumerate}

\[
\mathrm{M} ^ {\prime} \xrightarrow {u} \mathrm{M} \xrightarrow {v} \mathrm{M} ^ {\prime \prime} \longrightarrow 0
\]

an exact sequence of A-modules and \(n\) an integer \(>0\) . Let \(L = \operatorname{Coker}(T^n(u))\) so that there is an exact sequence

\[
\mathsf {T} ^ {n} (\mathbf {M} ^ {\prime}) \xrightarrow {\mathsf {T} ^ {n} (u)} \mathsf {T} ^ {n} (\mathbf {M}) \to \mathbf {L} \to 0.
\]

For every subset J of the set \(\{1, 2, \ldots, n\}\) , we write

\[
\mathrm{P} _ {J} = \mathrm{N} _ {1} \otimes \mathrm{N} _ {2} \otimes \dots \otimes \mathrm{N} _ {n}, \quad \text { where } \mathrm{N} _ {i} = \mathrm{M} ^ {\prime} \text { if } i \in \mathrm{J}, \mathrm{N} _ {i} = \mathrm{M} \text { if } i \notin \mathrm{J}
\]

\((\mathbf{P}_{\varnothing}\) is therefore equal to \(\mathsf{T}^n (\mathbf{M}))\) and

\[
\mathrm{Q} _ {\mathrm{J}} = \mathrm{N} _ {1} \otimes \mathrm{N} _ {2} \otimes \dots \otimes \mathrm{N} _ {n}, \quad \text { where } \mathrm{N} _ {i} = \mathrm{M} ^ {\prime} \text { if } i \in \mathrm{J}, \mathrm{N} _ {i} = \mathrm{M} ^ {\prime \prime} \text { if } i \notin \mathrm{J}
\]

\((Q_{\varnothing}\) is therefore equal to \(T^{n}(M^{\prime \prime})\) .) For every integer \(i\) such that \(0 \leqslant i \leqslant n\) , let \(P^{(i)}\) denote the submodule of \(T^{n}(M)\) generated by the union of the images of the \(\mathbf{P}_{\mathbf{J}}\) for all the subsets \(\mathbf{J}\) such that \(\operatorname{Card}(\mathbf{J}) = i\) and let \(\mathbf{L}^{(i)}\) be the canonical image of \(\mathbf{P}^{(i)}\) in \(\mathbf{L}\) .

\begin{enumerate}
\def\labelenumi{(\alph{enumi})}
\tightlist
\item
  Show that there exists a canonical surjective homomorphism
\end{enumerate}

\[
\bigoplus_ {\operatorname{Card} (J) = i} Q _ {J} \rightarrow L ^ {(i)} / L ^ {(i + 1)}
\]

(cf.~II, § 5, no. 6, Proposition 6).

\begin{enumerate}
\def\labelenumi{(\alph{enumi})}
\setcounter{enumi}{1}
\tightlist
\item
  If the sequence \(0 \to \mathbf{M}' \xrightarrow{u} \mathbf{M} \xrightarrow{v} \mathbf{M}'' \to 0\) is exact and splits, show that the homomorphism defined in (a) is bijective.
\end{enumerate}

